\section{Audit Protocol}

\subsection{Data}

The audit uses daily Yahoo Finance data from 2015-01-02 to 2026-04-30 after download cleaning. The cross-market universe contains 165 successful dataset-series memberships across five configured lists: Dow 30, Nasdaq large-technology stocks, an S\&P 100 subset, Hang Seng large-cap stocks, and a global ETF basket. Because some U.S. tickers appear in more than one configured list, these memberships correspond to 133 unique ticker symbols after cross-index overlaps. One configured Hong Kong ticker failed download and is excluded from the analyzed universe.

Two price modes are evaluated. The raw-close mode uses the downloaded OHLCV fields directly. The adjusted-close mode adjusts OHLC prices by the adjusted-close to close ratio, preserving the same volume field. Horizons are 1, 5, and 10 trading days. The target is the future absolute closing price level, not return, direction, volatility, or trading profit.

\input{tables/crossmarket_core_tables}

\subsection{Splitting and Baselines}

Each asset is split chronologically into 72\% training, 8\% validation, and 20\% testing. Feature scalers are fit only on the training segment and then applied to validation and test segments. Hyperparameter and checkpoint selection use validation loss only. Test labels are never shuffled.

The primary gate is latest-close persistence: at each test origin, the forecast equals the close at that origin. Additional no-training baselines, including drift-adjusted persistence and rolling means, are computed as diagnostics, but latest-close persistence is the mandatory minimum usefulness gate because it is simple, transparent, and matched to the absolute-price target.

For assets with the full downloaded calendar, the nominal data span is 2015-01-02 to 2026-04-30, with the chronological split falling approximately in early 2023 for the train/validation boundary and late 2023 for the validation/test boundary. Exact split rows vary slightly by market calendar and indicator availability and are stored in the evidence package. The adjusted-close mode multiplies same-day OHLC prices by the same-day adjusted-close to close ratio; it does not use future adjustment information beyond what is present in the downloaded daily record.

\subsection{Models and Implementation Scope}

The supervised audit includes LSTM, GRU, DLinear, scaled-down PatchTST-style, scaled-down iTransformer-style, and scaled-down TimeMixer-style implementations. The word ``style'' is important: these are compact audit implementations, not official full-capacity reproductions of the published architectures. They are included to stress-test the baseline-first claim across several architectural motifs, not to rank official model repositories. Therefore, the results should not be interpreted as a definitive performance assessment of the full PatchTST, iTransformer, or TimeMixer systems.

All supervised models use a direct forecasting setup: for each horizon, a separate target is constructed by shifting the close price by that horizon, rather than iterating a one-step model. The main cross-market runs use a 60-trading-day input window, hidden dimension 64, dropout 0.1 where applicable, AdamW with learning rate 0.001, early stopping on validation loss, and seeds 42, 123, and 777. LSTM and GRU use one recurrent layer. DLinear applies a linear temporal projection followed by a feature head. The PatchTST-style proxy uses patch length 10, stride 5, two transformer encoder layers, four attention heads, GELU activation, and feed-forward dimension four times the model dimension. The iTransformer-style proxy treats variables as tokens, embeds each variable's full temporal trace, and uses two transformer encoder layers with four heads. The TimeMixer-style proxy concatenates representations at scales 1, 2, and 4 and passes them through a two-layer GELU MLP.

Each supervised model is evaluated over all assets, two price modes, three horizons, two feature sets, and three random seeds, yielding 5,940 rows per supervised model. The feature sets are OHLCV and OHLCV plus RSI, MACD, and Bollinger-band indicators.

The zero-shot audit uses Chronos-T5-small with univariate close-price contexts of up to 512 observations and 20 generated samples per forecast. Because zero-shot generation is more expensive and does not have the same training loop as the supervised models, it is evaluated using a sparse sampled-origin protocol: up to 96 test origins are selected by evenly spacing indices across each test segment for each asset-price-horizon combination. This creates 990 Chronos rows. Chronos results are therefore reported separately as a foundation-model audit rather than merged as if they were identical to the supervised full rolling-origin protocol. With a full rolling-origin evaluation, Chronos pass rates and tail behavior could shift.

\subsection{Metrics and Statistical Tests}

For each row, the audit records RMSE, MAE, MAPE, and the ratio between model RMSE and latest-close RMSE. Ratios below 1.0 pass the persistence gate; ratios above 1.0 fail it. Model-level summaries report pass count, pass rate, median RMSE ratio, and bootstrap 95\% confidence intervals for the median. A one-sided Wilcoxon signed-rank test evaluates whether model RMSE ratios are greater than 1.0. Because rows share tickers, horizons, price modes, seeds, and model families, the inferential summaries are interpreted as paired diagnostic evidence under this audit design rather than as independent-sample claims over unrelated observations.
